\section{Introduction}
\label{sec:intro}

Earth's geomagnetic field (GMF), which typically ranges from \SI{25}{\micro\tesla} to \SI{65}{\micro\tesla} at the surface, has played a critical role in the evolution of the biosphere. Beyond its function in shielding the atmosphere from the solar wind and cosmic radiation, the GMF serves as an omnipresent environmental factor for life on Earth. Many species have developed sophisticated magnetoreception mechanisms for navigation and orientation, and the GMF is increasingly recognized as a foundational environmental pressure that influences fundamental biological processes at the molecular, cellular, and organismal levels \citep{Mo2012, Binhi2017}. 

In contrast, the Moon's magnetic environment is vastly different. While the Moon lacks a global dipole field today, localized crustal magnetic anomalies remain as remnants of an early lunar dynamo \citep{Wieczorek2006}. These crustal fields exhibit extreme heterogeneity across the lunar surface, with field strengths typically ranging from near zero to roughly \SI{20}{\nano\tesla} when measured at an altitude of \SI{30}{\kilo\meter} \citep{Tsunakawa2015, Hood2020}. Although significantly weaker than Earth's GMF, these localized fields can interact with the solar wind, occasionally forming ``mini-magnetospheres'' that may alter the local plasma and radiation environment \citep{Mitchell2008}.

From a biological perspective, any environment with a magnetic field substantially weaker than the GMF (typically $< \SI{5}{\micro\tesla}$) is classified as a hypomagnetic field (HMF). The space environment, including the lunar surface and interplanetary space, is inherently hypomagnetic. Understanding the biological effects of HMFs is increasingly recognized as a critical concern for deep space exploration. Prolonged exposure to HMF conditions in Earth-based simulations has been shown to induce a range of physiological and developmental effects, including altered plant growth, changes in microbiome community structure, and potential impacts on central nervous system function and DNA repair mechanisms \citep{Maffei2014, Mo2014}. 

As humanity prepares for sustained lunar surface operations through the Artemis program, comprehensive risk assessments for biological systems are essential. The Artemis III mission aims to land astronauts near the lunar South Pole \citep{NASA2022}, an area with unique illumination, volatile deposits, and magnetic characteristics. While the radiation and gravity environments of the Moon have been extensively studied, the extreme heterogeneity of the lunar magnetic field and its potential implications for biological experiments and human health remain poorly characterized.

This study aims to address this knowledge gap by evaluating the magnetic environments at candidate landing sites for past, present, and future lunar missions. We analyze crustal magnetic field data to characterize the field magnitudes at specific Artemis III candidate sites, as well as historic Apollo and Chang'e locations. We then synthesize these physical findings with a systematic review of the biological literature to assess the potential implications of lunar magnetic field heterogeneity. Our results provide a critical foundation for biological risk assessment and experimental design for future lunar surface missions.
